\section{One composite small-field step on the same analytic ball}
\label{cdm-section}

We combine Sections~\ref{cdj-section} and \ref{cdc-section} with the
actual composite constraint, Gaussian and integral of Notebook~27.
All statements are first finite-regulator statements on compatible
periodic tori. Constants are uniform in the preceding homogeneous
depth $M=\ell^k$ and volume, after the primitive odd ratio $\ell$,
macro ratio $R=\ell^s$ and auxiliary grouping $J=\ell^j$ are fixed.
The theorem concerns a local holomorphic small-field map. The measure
identification required beyond its original germ is stated explicitly.

Set $d=4$ and use the fixed trace-orthonormal real color coordinates.
The input activity graph consists of $J$-blocks of the current
unit lattice. The output graph consists of $J$-blocks of the next
lattice, equivalently $JR$-blocks of the current lattice, and is
also the Gaussian integration graph. Thus input and output use
the same grouping $J$ in their respective current units. Projection
$\pi_R$ has fiber $R^d$, preserves connected images, and satisfies
\begin{equation}
 |\pi_RH|\le 2^d+\frac{2^d}{k_R}|H|,
 \qquad k_R=1+\left\lfloor\frac{R-3}{4}\right\rfloor.
 \label{cdm-compression}
\end{equation}
The distinct lifted block interfaces are retained on the chosen
cofinal regulator family.

\subsection{Statement and exact field-dependent factors}
\label{cdm-statement-section}

Write the inherited total functional as
\begin{equation}
 E(a)=\sum_H E_H(a|_H),\qquad
 \norm{E}_{a_0;b}:=
 \sup_i\sum_{H\ni i}e^{a_0|H|}\norm{E_H}_b
 \le\epsilon\le\bar\epsilon,\qquad a_0=\frac1{16},
 \label{cdm-input}
\end{equation}
where the carriers are connected and each coefficient is holomorphic
on the scalar-coordinate product ball of radius $b$. The fixed
$\bar\epsilon<\infty$ bounds the input class. Assume local gauge
invariance of the \emph{total} $E$ on its original common domain.
If $E$ is a logarithm, its compatible holomorphic branch is part
of the input; coefficientwise gauge invariance is not assumed.

\begin{theorem}[A common-domain composite step]
\label{cdm-theorem}
There are finite choices $R,J$, positive $b,g_0$, and retained
link/color-ball thresholds $T=T(g)\ge1$ such that, for
$0<g<g_0$, the chosen small-field macrostep has a holomorphic
continuation to the entire next-angle ball $\norm{w}_\infty<b$.
With $W=e^{iw}$ it has the factorization
\begin{equation}
 \rho_{\mathrm{next}}(W)=
 N_{M,R,g}\exp\left[-g^{-2}S_{MR}(W)+E_{\mathrm{next}}(w)\right].
 \label{cdm-factorization}
\end{equation}
Here $N_{M,R,g}$ is field independent, and $S_{MR}$ is the
homogeneous ordinary classical minimum at total depth $MR$,
after the $d=4$ physical rescaling. No coupling renormalization
has yet been performed.

The output has connected strict $w$-activities with the same
angle radius $b$ and cardinality exponent $a_0$. In the notation
constructed below,
\begin{align}
 \norm{E_{\mathrm{next}}}_{a_0;b}
 &\le C_T\epsilon+(C_\Delta+C_{\det})b
       +2S_1\left[e^3\tau(T)+e^{v+q}(v+q)\right],
 \label{cdm-output-bound}\\
 v&\le C_{\mathrm{base}}(gT^3+gT+g^2T^2),\qquad
 q\le16C_TC_R\frac{gT}{b}\epsilon .
 \label{cdm-small-factors}
\end{align}
All constants are uniform in $M$ and compatible volume. The
Gaussian cutoff integral is nonzero on this whole ball and has
the analytic logarithm used in \eqref{cdm-factorization}.
The total correction satisfies the local gauge Ward identity
throughout the ball, and is invariant along gauge paths that
remain in it.

These are statements about the continuation of the actual local
small-field germ. They identify the exact chosen macro integral
throughout the ball when a gauge-orbit current density with the
stated pre-gauge representation is available. They do not identify
all quantum histories or the full Wilson pushforward merely from
holomorphic uniqueness.
\end{theorem}

The following construction proves the theorem and specifies its
constants and domains. The estimate is not an amplitude
contraction: $C_T$ can be much larger than one.

Use the actual composite ordinary map
$F_R=\mathcal M_\ell^s$ and the independent direct-$R$ contour
gauge. The physical backgrounds of Notebook~27 satisfy
\begin{equation}
 V_p(W)=\mathcal M_\ell^k U_{MR}(W),\qquad
 V^\star(W)=V_p(W)^{u(W)},\qquad P_R(V^\star)=1,
 \qquad S_M(V^\star)=S_M(V_p)=S_{MR}(W).
 \label{cdm-background}
\end{equation}
The last identity follows from the two actual minimum comparisons.
Current spacing is one and the next spacing is $R$; the
coupling-free action has no extra rescaling factor in dimension
four.

The retained chart preserves its nonpivot coordinates exactly.
The integration field is $e^{i\Phi_{V^\star}(g\xi)}V^\star$ and
the window is
$\prod_{b\ {\rm retained}}\mathbf1_{\{|\xi(b)|\le T\}}$.
The coupling stays inside $\Phi$. Set
\begin{align}
 A(w)&=
 D_z^2\left[
 S_M(e^{i\Phi_{V^\star}(z)}V^\star)
 +\mathcal G_R(e^{i\Phi_{V^\star}(z)}V^\star)\right]_{z=0},
 \notag\\
 E_{\mathrm{pre}}(w)&=E(a_p(w)),\qquad a_p=\frac1i\log V_p,
 \notag\\
 Z(w)&=\mathbb E_{A(w)^{-1}}\left[
 \chi_{\mathrm{window}}(\xi)
 e^{V_{\mathrm{base}}(w,\xi)+Q_{\mathrm{inh}}(w,\xi)}
 \right].
 \label{cdm-integral-parts}
\end{align}
The base exponent is the centered classical/gauge remainder,
the ambient pivot-straightening logarithms, and the Haar
logarithms. The inherited difference is exactly
\begin{align}
 Q_{\mathrm{inh}}(w,\xi)&=E(f_1(w,\xi))-E(f_0(w)),\notag\\
 f_\lambda(w,\xi)(b)&=
 \frac1i\log\left[
 e^{\,i\operatorname{Ad}_{u(w;b_-)^{-1}}
                \Phi_{V^\star(w)}(\lambda g\xi)(b)}
 V_p(w)(b)\right],\qquad f_0=a_p.
 \label{cdm-pre-gauge-insertion}
\end{align}
This group identity changes no integration variable $\xi$.

For clarity, let $n_e$ count positive current bonds, $n_c$ final
macro bonds, $n_s$ noncenter current sites, and
$n_z=\dim\mathfrak g\,(n_e-n_c)$. In trace-orthonormal exponential
coordinates write $d(e^{iX})=h_0j(X)dX$, $j(0)=1$, and put
\[
 c(g)=\int_G
 \exp[-(1-\operatorname{Re}\tau(h))/g^2]\,dh.
\]
With $\mathcal D_R(V)$ the actual pivot-delta determinant and
$A_0=A_{M,R,0}$ the flat retained precision, the field-independent
normalization before any additional vacuum subtraction is
\begin{equation}
 N_{M,R,g}=
 c(g)^{-n_s}h_0^{\,n_e-n_c}g^{n_z}(2\pi)^{n_z/2}
       \mathcal D_R(1)\det(A_0)^{-1/2},
 \label{cdm-normalization}
\end{equation}
times any separately tracked old field-independent normalization.
The background-dependent parts of the pivot and Gaussian
determinants remain in $E_{\mathrm{next}}$.

On the common overlap this is precisely the local Haar/delta
integral. If an actual gauge-invariant current density on gauge
orbits has representation $e^{-S_M/g^2+E}$ at the pre-gauge
fields in \eqref{cdm-pre-gauge-insertion}, the same expression
is the exact selected sector pointwise throughout the ball.
Otherwise it specifies only the verified analytic continuation
of that sector's local germ. The measurable completion and
one-shot sector are those of Notebook~27, not an asserted
intersection of $s$ sequential primitive sectors. Additional
original indicators would have to be retained and estimated.

\subsection{An explicit order of parameter choices}
\label{cdm-parameters-section}

First fix the exponents
\begin{equation}
 a_0=\frac1{16},\quad \kappa_G=2,\quad m_G=3,\quad
 \gamma=8,\quad B=\frac{53}{8},\quad \delta=1,\quad
 \kappa_{\mathrm{tail}}=8.
 \label{cdm-exponents}
\end{equation}
The tuned Gaussian density factor below has logarithm at most
$5/8$, so \eqref{cdj-Gaussian-weight} holds, and
$\kappa_{\mathrm{tail}}\ge B+\delta$.

Choose a fixed pre-gauge cardinal reserve
$A_{\mathrm{in}}\ge d\kappa_{\mathrm{tail}}$ before $R$.
By \eqref{cdj-ap-family}, there are $C_A,\kappa_0,r_A>0$ fixed
at the primitive ratio, independently of $R,M$ and volume, with
\begin{equation}
 \sup_i\sum_Xe^{\kappa_0\tau(\{i\}\cup X)+A_{\mathrm{in}}|X|}
                  \norm{a_{p,i;X}}_r
       \le C_A r/R,\qquad 0<r\le r_A.
 \label{cdm-pregauge-family}
\end{equation}
The norm takes the maximum over the whole current-link fiber
of one $R$-block \emph{inside} the support sum.
Choose a sufficiently large power $R=\ell^s$ so that
\begin{equation}
 R\ge8C_A,\qquad
 \frac{5^d2^d}{k_R}B\le a_0/2.
 \label{cdm-R-choice}
\end{equation}
The composite chart and its local radii are now fixed. The actual
flat retained Gaussian has uniform bounds
\begin{equation}
 \gamma_RI\le A_0\le B_RI,\qquad
 \lambda_0=B_R^{-1},\quad \Lambda_0=\gamma_R^{-1},
 \label{cdm-spectral-bounds}
\end{equation}
and an exponentially local inverse. These constants may depend
on $R$ but not on the old depth $M$.

Next choose any further \emph{classical-only} reserves needed
for the fixed-$R$ chart and the two lemmas in
Section~\ref{cdc-section}. They may depend on $R$ and use
independent classical outer radii. They do not change
$A_{\mathrm{in}}$ or $C_A$ in
\eqref{cdm-pregauge-family}--\eqref{cdm-R-choice}.
The rooted stationary field is a fixed-$R$ local analytic
function of $a_p$; its support norm is $O_R(r)$ with any such
fixed reserve. Lemma~\ref{cdc-precision-lemma} therefore
provides strict symmetric precision pieces with row norm
$O_R(r)$.

Choose $J=\ell^j$ sufficiently large for the grouped covariance
and precision exponents, in particular
\begin{equation}
 \frac{32}{J}\le\kappa',\qquad A'\ge64,\qquad
 \frac{16}{J}\le\kappa_0/2,
 \label{cdm-J-choice}
\end{equation}
with the available $\kappa',A'$ in
\eqref{cdc-precision-reserve}. Set
$\ell_{\mathrm{tail}}=c\log J$ plus a fixed offset, with
$c>4d/\kappa_0$, so that eventually
$\ell_{\mathrm{tail}}+2\le J$. The actual joint family of
Section~\ref{cdj-grouping-section} then has
\begin{align}
 s_{\mathrm{tail}}
 &\le J^d(b/4)e^{-\kappa_0\ell_{\mathrm{tail}}/2},
 &m_{\mathrm{core}}&=b/2,\notag\\
 c_0&=5^d2^d,&c_1&=5^d2^d/k_R,\notag\\
 h_0^{\mathrm{halo}}&=2d,&F&=R^d,\qquad
 D(J)=d\,\dim\mathfrak g\,J^d .
 \label{cdm-tail-data}
\end{align}
Here $h_0^{\mathrm{halo}}$ is a graph radius, distinct from the
Haar constant in \eqref{cdm-normalization}. Increase $J$ and
the tail offset until, with $t=s_{\mathrm{tail}}/m_{\mathrm{core}}$,
\begin{equation}
 t\le\tfrac12,\qquad D(J)\log\frac1{1-t}\le a_0/2.
 \label{cdm-tail-budget}
\end{equation}
This is possible because its sufficient bound is a constant
times $J^{2d}e^{-\kappa_0\ell_{\mathrm{tail}}/2}\to0$.
The ratio $t$ in this sufficient estimate is independent of $b$.
The entire local $\xi$ collar and the core lie in the fixed
two-neighbor halo used in \eqref{cdj-actual-core}.

Only now choose the final field radius. The macro retained
dimension and tunable density parameters are
\begin{align}
 D_G&=J^d d(R^d-1)\dim\mathfrak g,\qquad
 \theta=\frac1{4D_G},\notag\\
 \mathcal K_\theta
 &=\left[\frac{1+\theta}{(1-\theta)^3}\right]^{D_G/2},
 &\log\mathcal K_\theta&\le\frac58,\qquad
 \rho_\theta=\frac{1+\theta}{1-\theta}.
 \label{cdm-density-parameters}
\end{align}
The uniform Gaussian bounds fix a positive common strictly
local threshold $\eta_*$. Shrink it if necessary so that
$\eta_*\le1$.

Choose $b>0$ below all classical, physical-branch and chart
radii. In particular require
\begin{equation}
 \frac{eN_0e^{2\kappa_0}C_A b}{R\rho_R}\le\frac12,
 \label{cdm-joint-domain}
\end{equation}
as in \eqref{cdj-degree-majorant}. The substituted macro
precision norm $\nu(b)=O_{R,J}(b)$ must meet every strictly
local support-algebra threshold. Include
\begin{equation}
 \nu(b)\le\gamma_R/2,\qquad
 \nu(b)\le\frac{\theta\gamma_R^2}{4B_R},
 \label{cdm-precision-small}
\end{equation}
as justified in
\eqref{cdc-covariance-comparison}--\eqref{cdc-relative-neighborhood}.
Also impose the tunable density-algebra conditions
$u_G\le\theta$ and
$w_G\le\theta/[(1+\theta)\Lambda_0]$, where $u_G,w_G$
are its continuous $O(\nu)$ support-algebra quantities, and
the smallness for the determinant chain expansion below.
All hold after one shrink of $b$ with $R,J$ fixed.
This selects a single input and output radius, rather than
dividing the inherited radius during substitution.

A bad macrocell for the actual link/color-ball window implies
that its full Euclidean retained norm exceeds $T$. Its valid
joint-tail constant is therefore
\begin{equation}
 \tau(T)=\mathcal K_\theta\,2^{D_G/2}
       \exp\left[-\frac{T^2}{4\rho_\theta\Lambda_0}\right].
 \label{cdm-tail-constant}
\end{equation}
Choose $T\ge1$ with $e^3\tau(T)\le\eta_*/4$.
This estimates the bad event without replacing the actual
good window by a Euclidean macro ball.

Finally require the coupling to be small enough that
\begin{equation}
 L_*=\frac{b}{8C_RgT}\ge2,\qquad L_*gT<\rho_R,
 \label{cdm-lambda-radius}
\end{equation}
and the interaction bounds below hold uniformly for
$\epsilon\le\bar\epsilon$. The second condition is an
additional fixed restriction on $b$. For instance $T(g)$
may exceed the fixed threshold and grow as a fixed multiple
of $\sqrt{\log(1/g)}$. Then $\tau(T)\to0$, $gT^3\to0$ and
$gT/b\to0$ for this fixed $b$.

\subsection{Strict coefficients are already in the next background}
\label{cdm-support-section}

Lemma~\ref{cdc-precision-lemma} composes the actual precision
pieces with the supported stationary field while keeping their
matrix rows and columns in the original carrier. Its
scalar-coordinate factor is paid by the independent classical
cardinality reserve. Lemma~\ref{cdc-scalar-lemma} keeps the
whole spectator $\xi$ carrier even at zero background Taylor
degree, joins each substituted background support to it, and
reroots the resulting connected carrier. The classical action
includes its pair-supported quadratic part; its actual chart
Hessian retains both constraint and ordinary-log curvature
terms and the full absolute gauge Hessian.

Thus \eqref{cdc-base-smallness} supplies, directly in $(w,\xi)$,
\begin{equation}
 \norm{V_{\mathrm{base}}}_B\le v,\qquad
 v=C_{\mathrm{base}}(gT^3+gT+g^2T^2).
 \label{cdm-base-bound}
\end{equation}
The independent-background definition subtracts constant,
linear and quadratic Taylor terms before composition. At the
physical branch the total linear term is zero; it is not
silently dropped away from that branch.

For the inherited term, the actual complete-fiber construction
\eqref{cdj-joint-family}--\eqref{cdj-full-family} yields
\begin{equation}
 \norm{f_\lambda-a_p}_{\kappa_0,A_{\mathrm{in}};b,T}
                 \le C_R|\lambda|gT.
 \label{cdm-actual-family}
\end{equation}
Every coefficient reads $\xi$ only in the fixed output-root
collar, independent of its remote $w$ support. On
\eqref{cdm-lambda-radius}, the whole absolute family has
amplitude at most $b/4$. Hence every truncated core has that
same bound without using cancellation.
Theorem~\ref{cdj-transport} and
\eqref{cdj-actual-increment} give canonical connected carriers
\begin{align}
 \sum_K U_{K,\lambda}&=E(f_\lambda),&
 \norm{U_\lambda}_B&\le C_T\epsilon,\notag\\
 C_T&=e^{B5^d2^d+2d}R^dS_1,&
 S_1&=\sum_{z\in\mathbb Z^d}e^{-|z|_1},\notag\\
 \norm{Q_{\mathrm{inh}}}_B
 &\le q:=\frac{C_T\epsilon}{L_*-1}
 \le16C_TC_R\frac{gT}{b}\epsilon .
 \label{cdm-inherited-bound}
\end{align}
The $\lambda=0$ family is termwise $\xi$ independent and
represents $E_{\mathrm{pre}}$. Empty background support terms
of the fluctuation difference remain interactions if they
depend on $\xi$. Banach Cauchy is applied with its contour
integral before the positive anchored sum, as in
\eqref{cdj-Cauchy-increment}.

All these coefficients are already functions of the next
background $w$ on the Gaussian graph. The Gaussian output
will need no second stationary substitution and no second
generic inherited-radius estimate.

\subsection{Nonvanishing and a bound for the inherited logarithm difference}
\label{cdm-nonvanishing-section}

Combine equal-carrier exponents. Their exact multiplier norm is
\begin{equation}
 \sup_z\sum_{K\ni z}e^{B|K|}
       \norm{e^{V_{\mathrm{base},K}+Q_K}-1}
       \le e^{v+q}(v+q).
 \label{cdm-multiplier-bound}
\end{equation}
Products on distinct overlapping carriers are handled by the
overlap theorem. Each multiplier is defined on its actual
carrier window and extended by zero outside it. This is exact
on the global retained window. Since
$B\ge4+\log\mathcal K_\theta+\kappa_G$ and connected $K$
has metric spanning-tree length $|K|-1$, this norm pays the
strict Gaussian multiplier weight.

Impose
\begin{equation}
 v\le\frac{\eta_*}{16e},\qquad
 q\le\frac{\eta_*}{32e},\qquad e^3\tau(T)\le\eta_*/4.
 \label{cdm-quarter-budgets}
\end{equation}
With $\eta_*\le1$, \eqref{cdm-multiplier-bound} is below the
interaction quarter budget. The precision conditions have
already been imposed. The tunable strictly local Gaussian
theorem therefore proves nonvanishing of $Z$ on the full
$b$-ball and an analytic connected logarithm with
\begin{align}
 a_G&=\min\{1/2,\kappa_G/4\}=\frac12,\notag\\
 \norm{\log Z}_{a_G;b}
 &\le2S_1\left[e^3\tau(T)+e^{v+q}(v+q)\right].
 \label{cdm-log-bound}
\end{align}
This is a nonvanishing theorem on a volume-independent
complex domain, not an inference from real positivity.

There is also a precise inherited-log difference estimate.
For $q>0$, replace $Q_{\mathrm{inh}}$ by
$\zeta Q_{\mathrm{inh}}$ and set
\begin{equation}
 L_Q=\frac{\eta_*}{16e q}\ge2.
 \label{cdm-log-interpolation-radius}
\end{equation}
On $|\zeta|<L_Q$, the total exponent norm is at most
$\eta_*/(8e)$; since $\eta_*\le1$, its multiplier norm is
at most $\eta_*/8$. Together with the bad-block budget,
the Gaussian bound is at most $S_1\eta_*$ on that disk.
The canonical connected-log coefficients are holomorphic
as an anchored Banach family. Cauchy therefore gives
\begin{align}
 \norm{\log Z-\log Z_{\mathrm{base}}}_{a_G;b}
 &\le\frac{S_1\eta_*}{L_Q-1}
 \le32eS_1q\notag\\
 &\le512eS_1C_TC_R\frac{gT}{b}\epsilon .
 \label{cdm-log-difference}
\end{align}
For $q=0$ the difference is zero. This interpolation keeps
the base exponent throughout; it does not omit a factor
$e^{V_{\mathrm{base},K}}$ when comparing equal-carrier
multipliers.

\subsection{External normalizations and the exact update}
\label{cdm-update-section}

The pivot and Gaussian determinant corrections are
\begin{equation}
 \Delta(w)=\log\frac{\mathcal D_R(V^\star(w))}
                         {\mathcal D_R(1)},\qquad
 Q_{\det}(w)=-\frac12\log\det[A(w)A_0^{-1}].
 \label{cdm-external-logs}
\end{equation}
The first is the sum of actual local pivot log determinants
composed with the supported stationary field. The scalar
composition lemma gives connected coefficients with any
required fixed reserve and norm at most $C_\Delta b$.

For the second, put $C_0=A_0^{-1}$ and
$\mathcal B(w)=A(w)-A_0$ and expand afresh on the macro graph:
\begin{equation}
 Q_{\det}(w)=\frac12\sum_{n\ge1}\frac{(-1)^n}{n}
                         \operatorname{Tr}[(C_0\mathcal B(w))^n].
 \label{cdm-det-series}
\end{equation}
Join every entire precision label and each covariance path
by deterministic trees and shortest paths. Their cyclic
union is connected and has cardinality at most the total
edge-length budget plus one. The weighted row norm at
$\gamma=8$ and a rerooting spare give
\begin{equation}
 \norm{Q_{\det}}_{1/2;b}
 \le\frac{D_Ge^{1/2}S_1}{2}
       \frac{K_8\nu_8}{1-K_8\nu_8}
 \le C_{\det}b,\qquad K_8\nu_8<1.
 \label{cdm-det-bound}
\end{equation}
Here $K_8$ bounds the macro baseline covariance and $\nu_8$
the strict precision perturbation at that weight, uniformly
in $M$ and volume. Their smallness was included in the
choice of $b$.

With the field-independent factor
\eqref{cdm-normalization}, the exact correction is
\begin{equation}
 E_{\mathrm{next}}=
 E_{\mathrm{pre}}+\Delta+Q_{\det}+\log Z.
 \label{cdm-update}
\end{equation}
Its first term has weight $B$, and the other terms have
weight at least $1/2$, all on the same product ball $b$.
They therefore have the required weight $a_0=1/16$.
Equations \eqref{cdm-inherited-bound},
\eqref{cdm-log-bound} and \eqref{cdm-det-bound} prove
\eqref{cdm-output-bound}. Subtracting each output
coefficient's value at $w=0$ costs at most a factor two
in the norm and changes only the field-independent
normalization.

\subsection{Gauge covariance of the sector and Ward continuation}
\label{cdm-gauge-section}

Take a small real coarse gauge $k$ and lift it to current
sites as a constant on each final $R$-block. A background
with $P_R=1$ keeps that property under this lift. The
composite constraint transforms from $W$ to $W^k$, and
the classical action is invariant on the common
comparison domain. Uniqueness of the nearby physical
constrained minimum thus gives
\begin{equation}
 V^\star(W^k)=
 V^\star(W)^{\,\text{block-constant lift of }k}.
 \label{cdm-background-covariance}
\end{equation}
The relative retained chart transforms by source-endpoint
adjoints: this follows from constraint covariance,
unchanged pivots and uniqueness for the prescribed
retained values.

For real compact gauges these adjoints are orthogonal
in the chosen color norm. They preserve each actual
link/color-ball cutoff and have unit absolute retained
Jacobian. The original Haar/delta formula, equivalently
its exact pulled-back factors, therefore gives equality
of the chosen small-field integrals at $W$ and $W^k$
on a neighborhood of the identity. The inherited
identity is an identity of the total functional with
its compatible branch. No coefficientwise gauge
invariance is used.

The extracted classical action in
\eqref{cdm-background} is also gauge invariant there.
Nonvanishing gives the analytic logarithm on the whole
product ball. Differentiating the real near-origin
identity gives the infinitesimal coarse Ward identity
for $E_{\mathrm{next}}$. Its coefficients are holomorphic
functions of $w$. Uniqueness on the real slice, followed
by the holomorphic identity theorem on the connected
finite-volume ball, extends this identity throughout
$\norm{w}_\infty<b$. Integrating it gives invariance
along gauge paths staying in this ball.

This does not compare arbitrary endpoints whose
connecting gauge path leaves the domain. It also does
not turn the continued germ into a globally identified
current density without the orbit representation
specified above. A field-independent vacuum
normalization leaves the Ward identity unchanged.

\subsection{Finite-step scope}
\label{cdm-scope}

The common radius is obtained by placing the inherited
argument at the actual pre-gauge representative and
expanding only its distant supported tail. All new
classical and measure coefficients are already in
$w$ before integration. The Gaussian output therefore
incurs no second stationary transport. This proves
the stated domain, strict-support and nonvanishing
interfaces of one chosen composite step, uniformly
in preceding homogeneous refinement and volume.

The estimate still contains $C_T\epsilon$, not a
contractive amplitude update. Apart from optional
vacuum normalization, no relevant or marginal
extraction has been performed. A repeated quantum
construction requires a stable extraction and coupling
scheme, amplitude control, large-field and
mixed-history estimates, compatible global measure
identities, and the observable and continuum limits.
None follows here from analyticity or the local
Ward identity alone.
